\section{The adapter contract and the three tutors}\label{sec:adapters}

\subsection{Transport and session}

The adapter is the WebSocket \emph{client} and the Core is the \emph{server} at \code{ws://127.0.0.1:8765/ws/adapter}. This suits hosts that can open outgoing connections but would struggle to listen (Blender's embedded Python, a Chrome service worker). Messages are JSON-RPC 2.0 in both directions on one connection.

\begin{enumerate}[leftmargin=1.4em]
  \item The adapter sends \code{wysp.hello \{manifest, token\}}. The Core checks the token, the manifest and the major contract version.
  \item The adapter sends \code{wysp.heartbeat} at least every 20 s. The Core does not ping; it disconnects an adapter that has been silent for 45 s.
  \item The Core sends \code{adapter.*} requests. The adapter sends \code{event.*} notifications.
\end{enumerate}

{\small
\begin{xltabular}{\textwidth}{@{}L{4.9cm}L{3.4cm}L{3.0cm}Y@{}}
\toprule
\textbf{Method} & \textbf{Result} & \textbf{Capability} & \textbf{Budget} \\
\midrule
\code{adapter.observe} & Observation & \code{observe} & 1 s \\
\code{adapter.get\_context} & Context ($\leq$ 8 KB) & \code{context} & 0.5 s \\
\code{adapter.get\_actions} & ActionSpec list & \code{actions} & 0.5 s \\
\code{adapter.execute} & ActionResult, undo token & \code{execute} & 5 s \\
\code{adapter.undo} & ActionResult & \code{undo} & 5 s \\
\code{adapter.highlight} & ok & \code{highlight} & 0.5 s \\
\code{adapter.clear\_highlight} (1.1) & \code{\{cleared\}} & \code{highlight} & 0.5 s \\
\code{adapter.verify} & Evidence list & \code{verify.native} & 2 s \\
\midrule
\code{event.state\_changed} & \multicolumn{3}{>{\raggedright\arraybackslash}X}{\code{\{state\_version, changed[], origin\}}; capability \code{events.state}} \\
\code{event.selection\_changed} & \multicolumn{3}{>{\raggedright\arraybackslash}X}{\code{\{focus[]\}}; capability \code{events.selection}} \\
\code{event.user\_action} & \multicolumn{3}{>{\raggedright\arraybackslash}X}{\code{\{kind, target, details, origin\}}; capability \code{events.user\_action}} \\
\code{event.navigation} & \multicolumn{3}{>{\raggedright\arraybackslash}X}{\code{\{url, title\}}; capability \code{events.navigation}} \\
\code{event.highlight\_dismissed} (1.1) & \multicolumn{3}{>{\raggedright\arraybackslash}X}{\code{\{ref, reason\}}: the user touched the highlighted thing} \\
\bottomrule
\end{xltabular}}

Parameters: \code{observe \{scope\}} with a level (summary, focus, full), an optional root ref and a depth; \code{execute \{call\}} with action id, params, request id and \code{dry\_run}; \code{undo \{token\}}; \code{highlight \{ref, style, duration\_ms\}}.

Capabilities are promises: the Core calls only what the manifest declares. \code{observe}, \code{context} and \code{actions} are required. The core of the contract is deliberately small. Most adapters implement \code{observe}, \code{get\_context}, \code{get\_actions}, \code{execute}, \code{undo} and \code{highlight}, and verification is done by the Core on fresh observations, so no adapter re-implements verification logic. \code{adapter.verify} exists for cases where a native check is more reliable.

\subsection{The observation model}

An observation is an \emph{entity graph}: \code{entities[]} (\code{ref}, \code{kind}, \code{label}, \code{parent}, \code{props}, \code{confidence}), \code{links[]}, \code{focus[]}, \code{truncated} and \code{counts}. The Core knows no kinds in advance. Each adapter declares its kinds in the manifest, and the contract kit checks that every kind used is declared.

\begin{lstlisting}[language={}]
{"observation_id": "o-118", "state_version": 42, "scope": {"level": "focus"},
 "entities": [
  {"ref": "blender://object/Pumpkin", "kind": "object", "label": "Pumpkin",
   "props": {"type": "MESH", "modifiers_count": 2}},
  {"ref": "blender://object/Pumpkin/modifier/Solidify", "kind": "modifier",
   "label": "Solidify", "parent": "blender://object/Pumpkin",
   "props": {"type": "SOLIDIFY", "index": 1, "thickness": 0.02}}],
 "links": [], "focus": ["blender://object/Pumpkin"], "truncated": false}
\end{lstlisting}

Limits: 500 entities and 256 KB per observation, 8 KB for the context, refs matching \code{scheme://...} up to 512 characters. When an adapter cuts an observation it sets \code{truncated} and reports true totals in \code{counts}, and the tutor is told so it does not claim a list is complete. Entities whose \code{props.sensitive} is true are counted but never named.

The \textbf{state digest} (SHA-256 of canonical JSON over entities and links, without volatile properties, floats rounded to 6 significant digits) defines ``the state is unchanged'': an invalid call must leave it equal, and undo must restore it. Each adapter lists its own volatile properties in the manifest.

\subsection{Actions, risk levels and postconditions}

An adapter publishes a \emph{closed, declarative} catalog. An action declares an id (\code{verb\_object}), a description written for the model, a JSON Schema for its parameters (\code{additionalProperties: false}), a risk level, a \code{reversible} flag consistent with the risk, optional preconditions, \emph{postconditions} (predicates with \code{\$\{params.*\}} placeholders), and optionally \code{guide.ui\_path}, the steps the user would follow in the interface. Risk levels and their policy are in \cref{sec:gate}. An action that evaluates arbitrary code (Python \code{exec}, JavaScript injection, a shell command, an arbitrary IDE action id) is forbidden, and the kit rejects ids that look like it.

Execution rules every adapter follows: validate parameters before touching anything; support \code{dry\_run} (return a preview, change nothing, return no undo token); return an \code{undo\_token} for reversible actions and make \code{adapter.undo} restore the same digest; tag changes caused by Wysp with \code{origin: wysp} so the Core never credits them to the learner. Errors use the standard JSON-RPC codes plus \code{-32010} to \code{-32019} (not supported, action not found, target not found, precondition failed, execution failed, undo unavailable, timeout, busy, incompatible version, unauthorized).

\subsection{Persistent highlights (contract 1.1)}

\code{adapter.highlight} points at an entity or at an interface element the adapter documents (for example \code{intellij://ui/label/Commit}). Version 1.1 added \code{duration\_ms: 0}, meaning ``until cleared'', the request \code{adapter.clear\_highlight} and the notification \code{event.highlight\_dismissed}. An adapter keeps at most one highlight, removes it when the user interacts with the target (ignoring caret moves and scrolls its own highlight caused) and when the connection drops. The Core ends a persistent highlight when the tutor clears it, replaces it, a guided step is verified, the adapter disconnects, or after 10 minutes. With an adapter that declares 1.0, the Core falls back to a timed highlight.

\subsection{Versioning and generated artefacts}

\code{contract\_version} is 1.1.0. Additive, optional changes bump the minor version; breaking changes bump the major version. The Core accepts any adapter with the same major version. The Pydantic models are the source of truth: \code{uv run python -m wysp\_core.contract.export} writes the JSON Schemas to \file{schemas/} (26 schemas and an index) and the TypeScript types to \file{sdk/ts/src/generated/}. CI fails if the committed files differ from a fresh export. Neither directory is edited by hand.

\subsection{How to write an adapter}

The guide \file{docs/\allowbreak adapter-guide.md} builds a complete adapter for a toy notes application (\file{sdk/python/\allowbreak examples/\allowbreak notes\_adapter.py}, tested against the Core's own models). In short:

\begin{enumerate}[leftmargin=1.4em]
  \item \textbf{Manifest}: id, integration level, ref scheme, capabilities, entity kinds, volatile properties, profile id.
  \item \textbf{Observation}: map the application state to entities with stable refs. The Python SDK (\code{AdapterRuntime}) gives you the client, the state version and the routing; you implement \code{observe} and \code{get\_context}. In TypeScript, \code{@wysp/adapter-sdk} provides the client and \code{defineAction}.
  \item \textbf{Actions}: one file per action with schema, risk, postconditions and UI path.
  \item \textbf{Profile}: skills, concepts, lessons and exercises in YAML.
  \item \textbf{Check}: \code{uv run wysp-contract --tutor <name>} (\cref{sec:kit}).
\end{enumerate}

\subsection{What is application-specific}

\begin{figure}[H]
\centering
\begin{tikzpicture}[fig]
  \node[zone=lilacdeep, minimum width=15.4cm, minimum height=1.2cm, align=center] (core) at (0,2.2) {\textbf{Shared, unchanged: Wysp Core, Panel, protocol, SDK contract}};
  \foreach \x/\t in {-5.1/Blender, 0/IntelliJ IDEA, 5.1/Chrome} {
    \node[mod, minimum width=4.6cm, minimum height=0.7cm] at (\x,0.8) {\textbf{\t}\\[-1pt]\tiny adapter (code)};
    \node[mod, draw=okgreen, fill=okgreen!6, minimum width=4.6cm, minimum height=0.7cm] at (\x,-0.2) {\tiny profile (YAML data)};
  }
  \draw[arr] (-5.1,1.15) -- (-5.1,1.6); \draw[arr] (0,1.15) -- (0,1.6); \draw[arr] (5.1,1.15) -- (5.1,1.6);
\end{tikzpicture}
\caption{Adding a tutor changes only the two lower rows.}\label{fig:appspecific}
\end{figure}

\begin{xltabular}{\textwidth}{@{}L{2.4cm}Y Y Y@{}}
\toprule
 & \textbf{Blender} & \textbf{IntelliJ IDEA} & \textbf{Chrome (optional)} \\
\midrule
Code & Python extension, \code{bpy}; Blender 5.2 LTS & Kotlin plugin on the IntelliJ Platform 2026.2; no third-party runtime dependency & TypeScript, Manifest V3, Chrome 116 or later \\
Integration & \code{native} & \code{native} & \code{interface} \\
Ref scheme & \code{blender://} & \code{intellij://} & \code{web://} \\
Entity kinds & scene, collection, object, modifier, material & project, module, directory, file, symbol, problem, run configuration & page, region, heading, form, control, dialog, table, list, chart, status \\
Threading rule & Never touch \code{bpy} off the main thread. The network thread queues requests, a \code{bpy.app.timers} pump (every 50 ms) runs them. & Read the model in non-blocking read actions on pooled threads. Change code only in write commands on the EDT. Swing only on the EDT. & The WebSocket lives in the service worker (heartbeat 20 s keeps it alive). DOM work happens in the content script. \\
Undo & Blender's own undo history gets a step after each change, so Ctrl+Z also works. & Each change is one IDE command (``Wysp: ...''), read back after writing, restored if wrong. Wysp's undo never overwrites newer typing. & Undo tokens are stored in \code{storage.session}. \\
Pointing & GPU overlays: pulsing box on objects, frame and callout on editors, regions and Properties tabs. & Editor markup, a glass-pane painter and balloons; label positions come from Swing. & A closed shadow-DOM overlay in lilac that respects reduced motion. \\
\bottomrule
\end{xltabular}

\subsection{Blender}

The adapter runs inside Blender as an extension (\file{blender\_manifest.toml}, id \code{wysp\_tutor}, Blender 5.2 or later, GPL-3.0-or-later as Blender add-ons require) with two bundled pure-Python wheels: the adapter SDK and \code{websocket-client}. A depsgraph handler only marks the scene dirty; the main-thread loop diffs per-entity hashes (at most every 250 ms) and emits one \code{event.state\_changed}. Around each action the adapter flushes the diff before and after, so what the action changed is tagged \code{wysp} and everything else \code{user}. Blender cannot capture audio, so its panel uses the Core's \code{/api/listen/*} and \code{/api/tts} and plays the audio with \code{aud}.

Inside Blender the interface is native: a sidebar panel with the chat, an approval card and a viewport banner, and a node-tree view of the board (\code{.Wysp Board}, removed before the file is saved). It speaks the same \code{/ws/ui} protocol as the browser Panel, which is how a different UI technology still counts as the same Panel protocol. Actions (10): \code{select\_object}, \code{show\_properties\_tab} (\code{ui\_state}); \code{set\_object\_transform}, \code{add\_modifier}, \code{set\_modifier\_property}, \code{move\_modifier}, \code{toggle\_modifier\_visibility}, \code{assign\_material}, \code{remove\_modifier} (\code{reversible}); \code{apply\_modifier} (\code{irreversible}). Profile: 3 lessons, 4 skills, 12 concepts.

\subsection{IntelliJ IDEA}

The plugin connects with \code{java.net.http.WebSocket} and a small JSON codec, and embeds the Core's Panel in a tool window through JCEF (when JCEF is missing, the tool window links to the browser). The pairing token travels in the URL fragment, so no server ever sees it. Symbols come from the IDE's structure view, and imports (Java, Kotlin, Python) give the dependency edges used by exploration. While the IDE indexes, symbols are marked pending and the adapter sends \code{event.state\_changed} when indexing ends.

Actions (13): \code{open\_file}, \code{go\_to\_symbol}, \code{show\_tool\_window} (\code{ui\_state}); \code{read\_code}, \code{search\_code}, \code{find\_usages}, \code{list\_files} (\code{read\_only}, with caps such as 150 lines or 50 matches per call); \code{rename\_symbol}, \code{replace\_lines}, \code{insert\_lines}, \code{create\_file}, \code{reformat\_code}, \code{optimize\_imports} (\code{reversible}). Nothing runs code, run configurations or arbitrary IDE action ids. Sensitive files (\code{.env*}, keys, certificates, files named like credentials or secrets) are never read or written, by file name. Code text reaches the Core only for the selection (1500 characters), the code at the caret (3000) and read-tool results (8000 per call). Profile: 2 lessons, 6 skills, 22 concepts and an \code{explore} section.

\subsection{Chrome}

The extension is optional and sits outside the main demo path. It is active only on allowlisted origins (default \code{http://localhost:5180}, matched exactly), requested as optional host permissions per origin. The user pastes the pairing token into the options page. The content script builds an accessibility-oriented snapshot (roles and accessible names through \code{dom-accessibility-api}) and refs of the form \code{web://host/region\textasciitilde{}name/role\textasciitilde{}name}, rebuilt from the live DOM at execution time; an ambiguous target fails with up to five candidates.

\textbf{Redaction.} Values of password, hidden, card, one-time-code and \code{data-wysp-sensitive} fields are replaced with \code{[redacted]}; texts are cut at 200 characters; page URLs are sent without query or fragment; cookies, storage and request headers are never sent. Actions (7): \code{scroll\_into\_view}, \code{focus\_element}, \code{follow\_link} (\code{ui\_state}, same allowlisted site only); \code{fill\_field}, \code{select\_option}, \code{set\_checked} (\code{reversible}); \code{click\_element} (\code{irreversible}, with a dry-run preview). Profile: 3 lessons, 3 skills, 13 concepts.

\begin{wherebox}
\file{docs/adapter-guide.md}, \file{.kiro/specs/adapter-contract/}, \file{core/src/wysp\_core/contract/}, \file{schemas/}, \file{sdk/python/}, \file{sdk/ts/}. Tutors: \file{apps/blender-tutor/} (\file{adapter/runtime.py}, \file{adapter/actions/}), \file{apps/intellij-tutor/plugin/src/main/kotlin/dev/wysp/intellij/}, \file{apps/web-tutor/extension/src/}. Each has a \file{README.md} with its action table and privacy notes.
\end{wherebox}



